%% ─────────────────────────────────────────────────────────────────────────────
%% Lettre de motivation — Talvin Ackbaraly × MBDA France
%% Stage : démonstrateur embarqué de détection d'objets par IA à dimension temporelle
%% Même charte que le CV (XCharter, un seul accent, filets fins)
%% Compile: latexmk -pdf lettre-motivation-mbda.tex
%% ─────────────────────────────────────────────────────────────────────────────
\documentclass[11pt, a4paper]{extarticle}

\usepackage[T1]{fontenc}
\usepackage[utf8]{inputenc}
\usepackage[french]{babel}
\usepackage[top=1.6cm, bottom=1.6cm, left=2.2cm, right=2.2cm]{geometry}
\usepackage{XCharter}
\usepackage[letterspace=80]{microtype}
\usepackage{ragged2e}
\usepackage{xcolor}
\usepackage[hidelinks]{hyperref}

\hypersetup{pdftitle={Lettre de motivation — Talvin Ackbaraly}, pdfauthor={Talvin Ackbaraly}}

%% ── Colour system: identique au CV ──────────────────────────────────────────
\definecolor{accent}{HTML}{2563A0}
\definecolor{ink}{HTML}{1A1A1A}
\definecolor{soft}{HTML}{6E6E6E}
\definecolor{hair}{HTML}{C9C9C9}

\hypersetup{colorlinks=true, urlcolor=accent, linkcolor=accent}
\AtBeginDocument{\color{ink}}

\pagestyle{empty}
\setlength{\parindent}{0pt}
\setlength{\JustifyingParindent}{0pt}
\setlength{\parskip}{8pt}
\linespread{1.08}
\microtypesetup{expansion=false}

\begin{document}

%% ── En-tête : expéditeur ─────────────────────────────────────────────────────
%% Coordonnées seulement : l'école et les liens sont déjà dans le CV
{\color{accent}\fontsize{22}{26}\selectfont\scshape\textls[40]{Talvin Ackbaraly}}\par\vspace{3pt}
{\setlength{\parskip}{0pt}
94270 Le Kremlin-Bicêtre\\
\href{tel:+33769389008}{\textcolor{ink}{07 69 38 90 08}}\\
\href{mailto:talvin.ackbaraly@icloud.com}{\textcolor{ink}{talvin.ackbaraly@icloud.com}}\par}
\vspace{-4pt}{\color{hair}\rule{\textwidth}{0.4pt}}\par

%% ── Destinataire et date ─────────────────────────────────────────────────────
\vspace{2pt}
%% Bloc calé contre la marge droite, texte aligné à gauche dans le bloc
\hfill\begin{tabular}[t]{@{}l@{}}
  \textbf{MBDA France}\\
  Direction Technique\\
  1 avenue Réaumur\\
  92350 Le Plessis-Robinson
\end{tabular}

Le \today

\vspace{2pt}
\textbf{Objet : candidature au stage de fin d'études « Démonstrateur embarqué de détection d'objets
par IA à dimension temporelle » — 6 mois à partir de février 2027}

\vspace{2pt}
Madame, Monsieur,


\justifying
Étudiant en dernière année du cycle ingénieur à l'EPITA, dans la majeure Image, je vous adresse
ma candidature pour le stage consacré à un démonstrateur embarqué de détection d'objets par IA
à dimension temporelle, au sein de la Direction Technique de MBDA France. Ce sujet réunit les
deux domaines qui m'intéressent le plus : la détection par deep learning et l'optimisation de
code sur du matériel réel. J'aimerais aller de l'état de l'art jusqu'à un démonstrateur qui
tourne sur calculateur.

La détection d'objets par IA est au cœur de mon projet de fin d'études, mené avec la
Bibliothèque nationale de France. J'y suis en charge de la détection : j'ai comparé plusieurs
architectures, réglé les hyperparamètres et évalué les modèles sur un jeu de validation, jusqu'à
94,7\,\% de mAP50 sur 10\,000 images annotées. C'est la démarche que décrit votre offre : faire
l'état de l'art, retenir les méthodes les plus prometteuses et les tester sur des cas concrets.
J'ai aussi implémenté un U-Net de A à Z en PyTorch, sans modèle pré-entraîné.

La dimension temporelle est ce qui m'attire le plus dans ce stage. J'ai réalisé un projet de
détection de mouvement sur des flux vidéo avec des méthodes classiques : je connais donc les
approches auxquelles il faudra se comparer. Exploiter plusieurs images successives devrait aider
à détecter des objets petits, peu contrastés ou partiellement masqués, là où un détecteur image
par image échoue. Pendant l'état de l'art, j'aimerais étudier la fusion d'informations entre
images et les architectures qui traitent directement des séquences.

Le portage sur calculateur COTS est l'autre point fort de ce sujet pour moi. Dans ce même projet
de détection de mouvement, j'ai porté le traitement du CPU au GPU en CUDA et C++, avec une
accélération de ×24, en analysant les goulots d'étranglement avec Nsight. Je connais donc la
démarche : mesurer, identifier ce qui limite les performances, optimiser, puis vérifier que les
résultats restent conformes. Mon expérience en ingénierie logicielle (stage chez Studiocanal,
freelance) m'a aussi habitué à écrire du code propre et reproductible. Enfin, je suis à l'aise pour rédiger un rapport scientifique, en
français comme en anglais (TOEIC 845), ce qui compte dans l'environnement international de MBDA.

Je suis disponible pour un stage de six mois à partir de février 2027 et serais heureux de vous
présenter plus en détail ma motivation lors d'un entretien.

Je vous prie d'agréer, Madame, Monsieur, l'expression de mes salutations distinguées.

\vspace{10pt}
\hfill Talvin Ackbaraly\hspace*{1cm}

\end{document}
